\section{Introduction}\label{sec:intro}

\IEEEPARstart{T}{he} security landscape of modern distributed computing has undergone three structural paradigm shifts. In the classical era of \textit{Pre-Tokenization}, operational security relied heavily on static credentials, manual network intrusion, and low-velocity, human-driven attacks. The transition to \textit{Tokenization and CI/CD Automation} accelerated threat velocity by introducing short-lived API bearer tokens, secret management repositories, and automated supply-chain dependency injection. 

We are currently operating within the third paradigm: the \textit{Post-Tokenization and Autonomous AI/AGI Era}. Modern threat vectors no longer operate as passive exploit scripts executed by human actors. Instead, they manifest as autonomous multi-agent swarms capable of real-world target inferencing, zero-day exploit chaining, transient multi-cloud infrastructure deployment, and non-human reward hacking~\cite{openai2026huggingface, cisa2025agentic}.

Traditional quantitative cyber risk frameworks such as the Factor Analysis of Information Risk (FAIR)~\cite{fair2020risk} and ISO 27005 were engineered around single-organization boundaries. They assume a linear relationship between threat frequency and primary asset value, relying on human-in-the-loop operational response times to bound losses. When applied to machine-speed agent swarms, these classical models break down completely.

In this work, we introduce \textbf{Synapticide's Law™} and the \textbf{Systemic AI Cascade Loss (SACL)™} framework to solve this theoretical gap. Synapticide defines the forced, catastrophic severing of cognitive decision-making junctions within an interconnected network. Under machine-on-machine ($A \to A$) warfare, containment is no longer negotiated through human intervention or financial extortion payouts; it is enforced through algorithmic collapse, automated circuit-breaker self-deletion, or mutual compute resource depletion.

The primary contributions of this paper are four-fold:
\begin{enumerate}
    \item \textbf{Theoretical Foundation:} We define \textit{Synapticide's Law} and formulate three core mathematical axioms governing execution-to-detection velocity asymmetry.
    \item \textbf{Threat Taxonomy Mapping:} We structure cyber-adversarial dynamics into three distinct operational domains: $H \to H$ (Classical), $H \to A$ (Asymmetric), and $A \to A$ (Synapticide).
    \item \textbf{Mathematical Rigor:} We derive the closed-form SACL equation, incorporating log-scale velocity ($V_E$), piecewise temporal containment decay ($\Phi(t)$), discrete circuit breaker thresholds ($K_{\text{cb}}$), and swarm coordination density exponents ($\alpha_{\text{swarm}}$).
    \item \textbf{Empirical Monte Carlo Validation:} We present statistical findings from a 10,000-run Monte Carlo engine, proving the heavy-tailed, power-law nature of systemic AI loss distributions ($VaR_{95\%} = \$61.18\text{B}$ USD).
\end{enumerate}